%%%%%%%%%%%%%%%%%%%%%%%%%%%%%%%%%%%%%%%%%%%%%%%%%%%%%%%%%%%%%%%%%%%%%%%%%%%%%%%
%  SÍLABO DE ECONOMÍA INTERNACIONAL — Plan de estudio integrado
%
%  Comercio internacional:
%    FE  — Feenstra, Advanced International Trade: Theory and Evidence
%          (Princeton University Press, 2004)
%    KO  — Krugman, Obstfeld y Melitz, Economía internacional: teoría y
%          política (Pearson, 10.ª ed., 2016)
%    HK  — Helpman y Krugman, Trade Policy and Market Structure (MIT, 1989)
%  Macroeconomía internacional:
%    SW  — Schmitt-Grohé, Uribe y Woodford, International Macroeconomics
%          (borrador de mayo de 2016)
%    US  — Uribe y Schmitt-Grohé, Open Economy Macroeconomics (borrador de
%          marzo de 2015; Princeton University Press, 2017)
%    VE  — Végh, Open Economy Macroeconomics in Developing Countries
%          (MIT Press, 2013)
%    GA  — Galí, Monetary Policy, Inflation, and the Business Cycle
%          (Princeton University Press, 2008)
%  Programas PUCP integrados: ECO208 (Economía Internacional I), 1ECO16
%    (Economía Internacional 2), 1ECO83 (Economía Internacional, maestría),
%    1ECO93 (Finanzas Internacionales), 2ECO03 (Macroeconomía Avanzada),
%    1ECO74 (Macroeconomía Intermedia) y ECO338 (Tópicos de Macroeconomía
%    Avanzada 2).
%
%  Compilar con:  pdflatex EconomiaInternacional.tex   (dos veces, por el índice)
%%%%%%%%%%%%%%%%%%%%%%%%%%%%%%%%%%%%%%%%%%%%%%%%%%%%%%%%%%%%%%%%%%%%%%%%%%%%%%%

\documentclass[11pt,a4paper]{article}

\usepackage[T1]{fontenc}
\usepackage[utf8]{inputenc}
\usepackage[spanish,es-noquoting,es-tabla]{babel}
\usepackage{lmodern}
\usepackage[scaled=0.92]{helvet}        % títulos en sans serif, al estilo McGraw-Hill
\usepackage[margin=2.4cm,headheight=15pt]{geometry}
\usepackage[table]{xcolor}
\usepackage{enumitem}
\usepackage{titlesec}
\usepackage{titletoc}
\usepackage{tabularx}
\usepackage{booktabs}
\usepackage{longtable}
\usepackage{amsmath}
\usepackage{amssymb}
\usepackage{fancyhdr}
\usepackage{tikz}
\usepackage[skins]{tcolorbox}
\usepackage{needspace}

%------------------------------------------------------------------ colores ---
% Gama de rojos y burdeos
\definecolor{tonoOscuro}{RGB}{82,4,17}              % oscuro
\definecolor{tonoProfundo}{RGB}{121,6,26}           % profundo
\definecolor{tonoPrincipal}{RGB}{160,8,34}          % principal
\definecolor{tonoA}{RGB}{179,31,25}                 % variante cálida
\definecolor{tonoB}{RGB}{179,15,69}                 % variante fría
\definecolor{tonoC}{RGB}{193,21,107}                % variante viva
\definecolor{tonoClaro}{RGB}{221,95,116}            % claro
\definecolor{tonoAcero}{RGB}{131,63,74}             % acero
\definecolor{tonoPizarra}{RGB}{134,86,94}           % pizarra
\definecolor{tonoGrafito}{RGB}{73,39,45}            % grafito (texto secundario)
\definecolor{tonoHielo}{RGB}{247,207,214}           % hielo
\definecolor{tonoPalido}{RGB}{251,233,236}          % pálido
\definecolor{tonoBruma}{RGB}{254,246,247}           % bruma

\colorlet{azulosc}{tonoOscuro}
\colorlet{grisclaro}{tonoPalido}
% Un tono distinto de la gama para cada libro
\definecolor{colorFE}{RGB}{138,15,81}
\definecolor{colorKO}{RGB}{184,20,87}
\definecolor{colorHK}{RGB}{161,18,59}
\definecolor{colorSW}{RGB}{138,15,36}
\definecolor{colorUS}{RGB}{184,20,28}
\definecolor{colorVE}{RGB}{161,29,18}
\definecolor{colorGA}{RGB}{138,40,15}
\colorlet{colorPUCP}{tonoPizarra}

\usepackage[colorlinks=true,linkcolor=tonoOscuro,urlcolor=tonoPrincipal]{hyperref}

%------------------------------------------------------------------ estilos ---
\titleformat{\section}[block]
  {\sffamily\LARGE\bfseries\color{tonoOscuro}}{\thesection.}{0.6em}
  {\textcolor{tonoPrincipal}{\rule[-0.1ex]{0.55em}{0.85em}}\hspace{0.45em}}
  [\vspace{-0.4ex}{\color{tonoPrincipal}\titlerule[1.6pt]}\vspace{1.2pt}{\color{tonoClaro}\titlerule[0.6pt]}]
\titlespacing*{\section}{0pt}{2.6ex plus 1ex minus .2ex}{1.6ex plus .2ex}

\titleformat{\subsection}
  {\sffamily\large\bfseries\color{tonoPrincipal}}{\thesubsection}{0.6em}
  {\textcolor{tonoClaro}{\rule[0.1ex]{0.45em}{0.45em}}\hspace{0.45em}}
\titlespacing*{\subsection}{0pt}{2.2ex plus .8ex minus .2ex}{1ex plus .2ex}

% Índice
\titlecontents{section}[0pt]
  {\addvspace{3pt}\sffamily\bfseries\color{tonoOscuro}}
  {}{\textcolor{tonoB}{\rule[0.15ex]{0.4em}{0.4em}}\hspace{0.6em}}
  {{\color{tonoClaro}\titlerule*[0.7em]{.}}\textcolor{tonoPrincipal}{\contentspage}}

% Cabecera y pie
\pagestyle{fancy}
\fancyhf{}
\lhead{\small\sffamily\bfseries\color{tonoOscuro}Sílabo de Economía Internacional}
\rhead{\footnotesize\sffamily\color{tonoPrincipal}Feenstra $\cdot$ Krugman et al. $\cdot$ Uribe--Schmitt-Grohé $\cdot$ Végh}
\cfoot{{\setlength{\fboxsep}{3pt}\colorbox{tonoPrincipal}{\makebox[1.8em]{\small\sffamily\bfseries\color{white}\thepage}}}}
\renewcommand{\headrule}{%
  {\color{tonoPrincipal}\hrule height 1.1pt width\headwidth}\vskip 1.2pt
  {\color{tonoClaro}\hrule height 0.4pt width\headwidth}\vskip -2.7pt}

% Listas
\setlist[itemize,1]{label=\textcolor{tonoPrincipal}{\scriptsize$\blacksquare$}}
\setlist[enumerate,1]{label=\textcolor{tonoPrincipal}{\sffamily\bfseries\arabic*.}}

% Tablas: cabecera oscura y filas alternas en tono pálido
\arrayrulecolor{tonoPrincipal}
\setlength{\aboverulesep}{0pt}
\setlength{\belowrulesep}{0pt}
\renewcommand{\arraystretch}{1.3}
\newcommand{\cab}[1]{\textcolor{white}{\sffamily\bfseries #1}}
\newcommand{\filasalternas}{\rowcolors{2}{white}{tonoPalido}}

% Encabezados de párrafo en color
\newcommand{\destacado}[1]{{\sffamily\bfseries\color{tonoPrincipal}#1}}

% Franja decorativa con la gama
\newcommand{\franjas}[1][3pt]{%
  \noindent\begin{tikzpicture}
    \foreach \c [count=\i from 0] in {tonoOscuro,tonoProfundo,tonoPrincipal,tonoA,tonoB,tonoC,tonoClaro,tonoHielo}
      \fill[\c] (\i*\linewidth/8,0) rectangle ++(\linewidth/8,#1);
  \end{tikzpicture}\par}

%---------------------------------------------------------------- comandos ----
% Etiqueta de libro: recuadro de color con la sigla
\newcommand{\etiqueta}[2]{%
  {\setlength{\fboxsep}{1.4pt}\colorbox{#1}{\sffamily\bfseries\color{white}\strut#2}}}
% Referencias abreviadas a cada libro
\newcommand{\FE}[1]{\etiqueta{colorFE}{FE}\nobreak\hspace{0.25em}\textcolor{colorFE}{#1}}
\newcommand{\KO}[1]{\etiqueta{colorKO}{KO}\nobreak\hspace{0.25em}\textcolor{colorKO}{#1}}
\newcommand{\HK}[1]{\etiqueta{colorHK}{HK}\nobreak\hspace{0.25em}\textcolor{colorHK}{#1}}
\newcommand{\SW}[1]{\etiqueta{colorSW}{SW}\nobreak\hspace{0.25em}\textcolor{colorSW}{#1}}
\newcommand{\US}[1]{\etiqueta{colorUS}{US}\nobreak\hspace{0.25em}\textcolor{colorUS}{#1}}
\newcommand{\VE}[1]{\etiqueta{colorVE}{VE}\nobreak\hspace{0.25em}\textcolor{colorVE}{#1}}
\newcommand{\GA}[1]{\etiqueta{colorGA}{GA}\nobreak\hspace{0.25em}\textcolor{colorGA}{#1}}

% Un tema: número en recuadro de color + título
\newcommand{\tema}[1]{\temaaux#1\relax}
\def\temaaux#1 #2\relax{%
  \par\addvspace{1.1em}\noindent
  {\sffamily\bfseries\large
   \hangindent=2.9em\hangafter=1
   \makebox[2.9em][l]{{\setlength{\fboxsep}{2.5pt}\colorbox{tonoPrincipal}{\color{white}#1}}}%
   {\color{tonoOscuro}#2}\par}%
  \nobreak\vspace{0.25em}\noindent{\color{tonoClaro}\rule{\linewidth}{0.6pt}}\par\nobreak\vspace{0.35em}}

% Línea de lecturas con marcador y etiqueta de color
\newcommand{\linealectura}[3]{%
  \par\nobreak\noindent\hspace{1.2em}%
  \begin{minipage}{0.92\textwidth}\small
    \textcolor{#1}{\rule[0.1ex]{0.45em}{0.45em}}\hspace{0.45em}%
    {\sffamily\bfseries\color{#1}#2}\enspace#3%
  \end{minipage}\par\vspace{0.3em}}
\newcommand{\lecturas}[1]{\linealectura{tonoPrincipal}{Lecturas:}{#1}}
\newcommand{\apoyo}[1]{\linealectura{tonoA}{Profundización:}{#1}}
\newtcolorbox{cajanota}{enhanced, sharp corners, boxrule=0pt, frame hidden,
  colback=tonoPalido, borderline west={2.5pt}{0pt}{tonoClaro},
  left=8pt, right=6pt, top=4pt, bottom=4pt, before skip=0pt, after skip=0pt,
  fontupper=\small\color{tonoGrafito}}
\newcommand{\nota}[1]{%
  \par\nobreak\noindent\hspace{1.2em}%
  \begin{minipage}{0.92\textwidth}\begin{cajanota}#1\end{cajanota}\end{minipage}\par\vspace{0.3em}}
\newcommand{\consultalib}[1]{\linealectura{tonoPizarra}{Consulta:}{{\color{tonoGrafito}#1}}}
% Artículos y textos tomados de los programas PUCP
\newcommand{\articulos}[1]{\linealectura{colorPUCP}{Artículos (PUCP):}{#1}}
\newcommand{\consulta}[1]{\linealectura{colorPUCP}{Consulta (PUCP):}{{\color{tonoGrafito}#1}}}

% Encabezado de bloque: banda oscura, franja de la gama y recuadro de presentación
\newif\ifrecienparte
\newcommand{\bloque}[2]{%
  \ifrecienparte\global\recienpartefalse\else\par\needspace{0.4\textheight}\fi
  \phantomsection\addcontentsline{toc}{section}{#1}%
  \begin{tcolorbox}[enhanced, sharp corners, boxrule=0pt, frame hidden,
      colback=tonoOscuro, borderline west={7pt}{0pt}{tonoClaro},
      left=14pt, right=10pt, top=10pt, bottom=10pt, before skip=0pt, after skip=0pt]
    {\raggedright\sffamily\Large\bfseries\color{white}#1\par}
  \end{tcolorbox}%
  \nointerlineskip\franjas[4pt]
  \vspace{0.8em}
  \begin{tcolorbox}[enhanced, sharp corners, boxrule=0pt, frame hidden,
      colback=tonoPalido, borderline west={3pt}{0pt}{tonoPrincipal},
      borderline south={0.6pt}{0pt}{tonoClaro},
      left=10pt, right=8pt, top=6pt, bottom=6pt, before skip=0pt, after skip=0pt,
      fontupper=\small\color{tonoGrafito}]
    #2
  \end{tcolorbox}%
  \par\medskip}

% Encabezado de parte: banda a todo el ancho en el tono principal
\newcommand{\parte}[1]{%
  \clearpage
  \phantomsection\addcontentsline{toc}{section}{\textbf{#1}}%
  \begin{tcolorbox}[enhanced, sharp corners, boxrule=0pt, frame hidden,
      colback=tonoPrincipal, halign=center, top=16pt, bottom=16pt,
      borderline north={3pt}{0pt}{tonoClaro}, borderline south={3pt}{0pt}{tonoClaro},
      before skip=0pt, after skip=0pt]
    {\sffamily\LARGE\bfseries\color{white}#1}
  \end{tcolorbox}\par\bigskip
  \global\recienpartetrue}

%================================================================== documento ==
\begin{document}

\begin{titlepage}
\begin{tikzpicture}[remember picture,overlay]
  % banda superior
  \fill[tonoOscuro] (current page.north west) rectangle ([yshift=-5.4cm]current page.north east);
  \fill[tonoClaro] (current page.north west) rectangle ([xshift=1.1cm,yshift=-5.4cm]current page.north west);
  % franja de la gama bajo la banda
  \foreach \c [count=\i from 0] in {tonoOscuro,tonoProfundo,tonoPrincipal,tonoA,tonoB,tonoC,tonoClaro,tonoHielo}
    \fill[\c] ([xshift=\i*\paperwidth/8,yshift=-5.4cm]current page.north west)
      rectangle ++(\paperwidth/8,-0.35cm);
  % franja inferior
  \foreach \c [count=\i from 0] in {tonoHielo,tonoClaro,tonoC,tonoB,tonoA,tonoPrincipal,tonoProfundo,tonoOscuro}
    \fill[\c] ([xshift=\i*\paperwidth/8]current page.south west)
      rectangle ++(\paperwidth/8,0.6cm);
\end{tikzpicture}%
\centering
\vspace*{0.5cm}
{\sffamily\Huge\bfseries\color{white} Sílabo de Economía Internacional\par}
\vspace{0.4cm}
{\sffamily\Large\color{tonoHielo} Plan de estudio integrado\par}
\vspace{1.9cm}
{\sffamily\large\bfseries\color{tonoPrincipal}\textcolor{tonoClaro}{\rule[0.5ex]{1.5cm}{0.6pt}}\quad Comercio internacional\quad\textcolor{tonoClaro}{\rule[0.5ex]{1.5cm}{0.6pt}}\par}
\vspace{0.3cm}
{\large\etiqueta{colorFE}{FE} \quad {\color{tonoOscuro}Robert C. Feenstra}\par}
{\normalsize \emph{\color{colorFE}Advanced International Trade: Theory and Evidence}\par}
{\small\color{tonoPizarra} Princeton University Press, 2004 --- 11 capítulos y 2 apéndices\par}
\vspace{0.2cm}
{\large\etiqueta{colorKO}{KO} \quad {\color{tonoOscuro}Paul R. Krugman, Maurice Obstfeld y Marc J. Melitz}\par}
{\normalsize \emph{\color{colorKO}Economía internacional: teoría y política}\par}
{\small\color{tonoPizarra} Pearson, 10.ª edición, 2016 --- 22 capítulos\par}
\vspace{0.2cm}
{\large\etiqueta{colorHK}{HK} \quad {\color{tonoOscuro}Elhanan Helpman y Paul R. Krugman}\par}
{\normalsize \emph{\color{colorHK}Trade Policy and Market Structure}\par}
{\small\color{tonoPizarra} The MIT Press, 1989 --- 9 capítulos\par}
\vspace{0.6cm}
{\sffamily\large\bfseries\color{tonoPrincipal}\textcolor{tonoClaro}{\rule[0.5ex]{1.5cm}{0.6pt}}\quad Macroeconomía internacional\quad\textcolor{tonoClaro}{\rule[0.5ex]{1.5cm}{0.6pt}}\par}
\vspace{0.3cm}
{\large\etiqueta{colorSW}{SW} \quad {\color{tonoOscuro}Stephanie Schmitt-Grohé, Martín Uribe y Michael Woodford}\par}
{\normalsize \emph{\color{colorSW}International Macroeconomics}\par}
{\small\color{tonoPizarra} Borrador de mayo de 2016 --- 13 capítulos\par}
\vspace{0.2cm}
{\large\etiqueta{colorUS}{US} \quad {\color{tonoOscuro}Martín Uribe y Stephanie Schmitt-Grohé}\par}
{\normalsize \emph{\color{colorUS}Open Economy Macroeconomics}\par}
{\small\color{tonoPizarra} Borrador de marzo de 2015 (Princeton University Press, 2017) --- 11 capítulos\par}
\vspace{0.2cm}
{\large\etiqueta{colorVE}{VE} \quad {\color{tonoOscuro}Carlos A. Végh}\par}
{\normalsize \emph{\color{colorVE}Open Economy Macroeconomics in Developing Countries}\par}
{\small\color{tonoPizarra} The MIT Press, 2013 --- 17 capítulos\par}
\vspace{0.5cm}
{\normalsize\color{colorPUCP} Con temas y lecturas de siete programas de la PUCP\par}
\vfill
{\setlength{\fboxsep}{6pt}\colorbox{tonoPrincipal}{\normalsize\sffamily\bfseries\color{white} Nueve bloques temáticos --- 6 meses}\par}
\vspace{0.2cm}
\end{titlepage}

\tableofcontents
\newpage

%------------------------------------------------------------------------------
\section*{Cómo usar este sílabo}
\addcontentsline{toc}{section}{Cómo usar este sílabo}

El plan está organizado en \textbf{nueve bloques temáticos}, agrupados en dos
partes. La \textbf{Parte I} (Bloques 1 a 4) es el comercio internacional: de la
ventaja comparativa ricardiana y las dotaciones de factores a la competencia
imperfecta, las empresas heterogéneas, la ecuación de gravedad y la política
comercial. La \textbf{Parte II} (Bloques 5 a 9) es la macroeconomía
internacional desde la perspectiva de una economía pequeña, abierta y
emergente: la cuenta corriente, el tipo de cambio real, los ciclos económicos,
la política cambiaria, los flujos de capital, el sobreendeudamiento y la deuda
soberana. Conviene cerrar cada bloque antes de abrir el siguiente.

\medskip
\noindent\destacado{Por qué un curso aparte.} El comercio internacional no aparece
en ningún otro sílabo del conjunto. La macroeconomía internacional, en cambio,
aparece en dos, pero acotada: el sílabo de macroeconomía la limita al
Bloque 3 (cuenta corriente intertemporal y tipo de cambio real con De
Gregorio), deja fuera los capítulos de economía abierta de Romer y de
Campante et al.\ y cita solo como consulta, sin horas, los de Acemoglu,
Walsh y Galí; el de economía financiera dedica su
Bloque 8 a las finanzas internacionales con Krugman et al. Este plan
desarrolla ambos frentes con los textos de posgrado.

\medskip
\noindent\destacado{Relación con los demás sílabos.} Los capítulos 14 a 16, 18 a 20
y 22 de Krugman et al.\ ya se leen en el Bloque 8 del sílabo de economía
financiera; aquí se citan como consulta para repasarlos. La Parte I supone la
teoría del consumidor, el equilibrio general y la competencia monopolística del
sílabo de microeconomía; la Parte II, el modelo de Ramsey (Bloque 4) y el
modelo neokeynesiano (Bloque 13) del sílabo de macroeconomía, y las
herramientas de los Bloques 9 a 13 del sílabo de economía matemática.

\medskip
\noindent\destacado{Criterio general de reparto} (en el orden en que se lee cada tema):
\begin{itemize}[leftmargin=1.4em,itemsep=2pt]
  \item \textcolor{colorKO}{\textbf{Krugman, Obstfeld y Melitz}} abre cada tema
        de la Parte I con la intuición, los gráficos y la evidencia. Es el
        texto de ECO208 y de 1ECO16. Se lee \emph{primero}.
  \item \textcolor{colorFE}{\textbf{Feenstra}} es el texto base de la Parte I:
        da la versión formal de cada modelo y, sobre todo, su contrastación
        empírica. Es el texto de 1ECO16 y del Módulo 1 de 1ECO83.
  \item \textcolor{colorHK}{\textbf{Helpman y Krugman}} es la referencia de la
        política comercial con competencia imperfecta (Bloque 4).
  \item \textcolor{colorSW}{\textbf{Schmitt-Grohé, Uribe y Woodford}} es el
        texto de entrada de la Parte II: cuentas externas, cuenta corriente
        intertemporal, tipo de cambio real, deuda externa y dinero, con
        modelos de dos períodos y mucha evidencia.
  \item \textcolor{colorUS}{\textbf{Uribe y Schmitt-Grohé}} es el texto
        principal de la Parte II en su nivel de posgrado: ciclos en economías
        emergentes, términos de intercambio, rigidez salarial y política
        cambiaria, controles de capital, sobreendeudamiento y deuda soberana.
  \item \textcolor{colorVE}{\textbf{Végh}} aporta la tradición de la
        macroeconomía de los países en desarrollo: enfoque monetario de la
        balanza de pagos, estabilización de la inflación alta, flujos de
        capital, dolarización y crisis. Se lee como profundización.
\end{itemize}

\medskip
\noindent\destacado{Qué entra en las horas.} Las líneas \emph{Lecturas},
\emph{Profundización} y \emph{Artículos (PUCP)} entran en el cómputo de horas.
Las líneas \emph{Consulta} y \emph{Consulta (PUCP)} remiten a libros de la
carpeta y a referencias de los programas que no entran (Anexo E).

\medskip
\noindent\destacado{Material excluido.} Quedan fuera las finanzas internacionales
que ya cubre el sílabo de economía financiera, los métodos numéricos de
solución de modelos y los temas de coyuntura de los programas (fragmentación
geoeconómica, guerra comercial, informes de organismos), salvo los que
cuentan con un texto analítico. El detalle está en el Anexo B.

\medskip
\noindent\destacado{Notación.} \FE{n}, \KO{n}, \VE{n} y \GA{n} remiten al
capítulo $n$ de Feenstra, Krugman et al., Végh y Galí: ninguno de ellos
numera sus secciones en el índice, así que se citan por capítulo. \HK{n.m},
\SW{n.m} y \US{n.m} remiten a la sección $m$ del capítulo $n$ de Helpman y
Krugman, Schmitt-Grohé et al.\ y Uribe y Schmitt-Grohé. \emph{Consulta}
remite a libros de la carpeta que no entran en el cómputo de horas
(Anexo E).

%------------------------------------------------------------------------------
\newpage
\section*{Cronograma por semanas}
\addcontentsline{toc}{section}{Cronograma por semanas}

Calendario del plan: 260 horas distribuidas en \textbf{26 semanas de
contenido}, que equivalen a \textbf{6 meses de calendario}, a un ritmo de diez
horas semanales.

\medskip
\noindent Todas las semanas son de materia nueva y no hay semanas de colchón.
La columna de horas es la que gobierna el reparto. Si hay semanas perdidas por
feriados, viajes o imprevistos, el ajuste correcto es alargar el calendario,
no comprimir los bloques.

\medskip
{\filasalternas
\noindent\begin{tabularx}{\textwidth}{lXcc}
\toprule
\rowcolor{tonoOscuro}\cab{Semanas} & \cab{Contenido} & \cab{Sem.} & \cab{Horas} \\
\midrule
\rowcolor{tonoA}\multicolumn{4}{l}{\cab{\textit{Parte I --- Comercio internacional}}} \\
1--2   & Bloque 1 --- Comercio mundial y ventaja comparativa ricardiana & 2 & 20 \\
3--5   & Bloque 2 --- Dotaciones de factores y distribución del ingreso & 3 & 30 \\
6--8   & Bloque 3 --- Rendimientos crecientes, empresas y gravedad & 3 & 30 \\
9--11  & Bloque 4 --- Política comercial & 3 & 30 \\
\midrule
\rowcolor{tonoA}\multicolumn{4}{l}{\cab{\textit{Parte II --- Macroeconomía internacional}}} \\
\rowcolor{tonoHielo}\textbf{12--15} & \textbf{Bloque 5 --- La cuenta corriente} & \textbf{4} & \textbf{40} \\
16--17 & Bloque 6 --- Integración financiera y tipo de cambio real & 2 & 20 \\
18--20 & Bloque 7 --- Ciclos económicos en economías abiertas y emergentes & 3 & 30 \\
21--23 & Bloque 8 --- Dinero, tipo de cambio y política cambiaria & 3 & 30 \\
24--26 & Bloque 9 --- Flujos de capital, crisis y deuda soberana & 3 & 30 \\
\midrule
\rowcolor{tonoHielo}\multicolumn{2}{l}{\textbf{Total}} & \textbf{26} & \textbf{260} \\
\bottomrule
\end{tabularx}
}

\medskip
\noindent\destacado{Ritmo supuesto.} Diez horas por semana, todas de contenido
nuevo: por ejemplo, cuatro sesiones de dos horas entre semana y una sesión de
dos horas el sábado para ejercicios y datos.

\medskip
\noindent\destacado{El reparto entre las dos partes.} El comercio ocupa 11
semanas y la macroeconomía internacional, 15. La diferencia se explica porque
la Parte II reúne tres textos de posgrado y porque la Parte I se apoya en la
microeconomía ya estudiada.

\medskip
\noindent\destacado{El tramo largo.} El Bloque 7 (semanas 18--20) es el núcleo
cuantitativo del plan: el modelo de ciclos reales de economía pequeña y
abierta exige la log-linealización y los sistemas dinámicos de los sílabos de
macroeconomía (Bloque 14) y de economía matemática (Bloque 10).

\medskip
\noindent\destacado{Hitos de control.} Conviene verificar el avance al terminar las
semanas 11 y 20. Si en alguno de esos cortes se acumula más de una semana de
retraso, se alarga el calendario y no se comprime el bloque siguiente.

\subsection*{Detalle semana a semana}

{\small\filasalternas
\renewcommand{\arraystretch}{1.1}\noindent\begin{tabularx}{\textwidth}{clXc}
\toprule
\rowcolor{tonoOscuro}\cab{Sem.} & \cab{Bloque} & \cab{Temas} & \cab{Horas} \\
\midrule
1  & 1 & 1.1 El comercio mundial: una visión general; 1.2 El modelo ricardiano & 10 \\
2  & 1 & 1.3 El modelo ricardiano con un continuo de bienes y su evidencia & 10 \\
\midrule
3  & 2 & 2.1 El modelo de factores específicos; 2.2 El modelo de Heckscher--Ohlin y sus teoremas & 10 \\
4  & 2 & 2.3 Muchos bienes y factores: el modelo de Heckscher--Ohlin--Vanek; 2.4 La evidencia empírica & 10 \\
5  & 2 & 2.5 El modelo estándar, los términos de intercambio y los auges de materias primas & 10 \\
\midrule
6  & 3 & 3.1 Economías externas de escala y localización; 3.2 Competencia monopolística y comercio intraindustrial & 10 \\
7  & 3 & 3.3 La ecuación de gravedad; 3.4 Empresas heterogéneas: el modelo de Melitz & 10 \\
8  & 3 & 3.5 Comercio de insumos intermedios, deslocalización y salarios; 3.6 Empresas multinacionales & 10 \\
\midrule
9  & 4 & 4.1 Instrumentos de política comercial y protección efectiva; 4.2 Política comercial con competencia imperfecta & 10 \\
10 & 4 & 4.3 Economía política de la política comercial; 4.4 Acuerdos comerciales regionales & 10 \\
11 & 4 & 4.5 Comercio, desarrollo y crecimiento; 4.6 Comercio y desigualdad: evidencia y el caso peruano & 10 \\
\midrule
12 & 5 & 5.1 Balanza de pagos y posición de inversión internacional; 5.2 Sostenibilidad de la cuenta corriente & 10 \\
13 & 5 & 5.3 La cuenta corriente en una economía de dotación; 5.4 Inversión y cuenta corriente & 10 \\
14 & 5 & 5.5 Incertidumbre y prueba del enfoque intertemporal; 5.6 Ajuste externo en economías pequeñas y grandes & 10 \\
15 & 5 & 5.7 Déficits gemelos y equivalencia ricardiana; 5.8 Imperfecciones del mercado de capitales & 10 \\
\midrule
16 & 6 & 6.1 Movilidad del capital y paridades de tasas de interés; 6.2 Determinantes del tipo de cambio real & 10 \\
17 & 6 & 6.3 Bienes no transables y precios intertemporales; 6.4 El modelo TNT y las paradas súbitas & 10 \\
\midrule
18 & 7 & 7.1 Hechos estilizados del ciclo en el mundo; 7.2 El modelo RBC de economía pequeña y abierta & 10 \\
19 & 7 & 7.3 Ciclos en economías emergentes; 7.4 Choques de tasa de interés & 10 \\
20 & 7 & 7.5 Términos de intercambio; 7.6 Ciclos internacionales y los enigmas de la macroeconomía abierta & 10 \\
\midrule
21 & 8 & 8.1 Dinero, inflación y tipo de cambio nominal; 8.2 Producción a corto plazo y áreas monetarias óptimas & 10 \\
22 & 8 & 8.3 Rigidez salarial, desempleo y política cambiaria óptima; 8.4 El modelo neokeynesiano de economía abierta & 10 \\
23 & 8 & 8.5 Regímenes cambiarios y miedo a flotar; 8.6 Estabilización de la inflación alta & 10 \\
\midrule
24 & 9 & 9.1 Flujos de capital y paradas súbitas; 9.2 Controles de capital & 10 \\
25 & 9 & 9.3 Sobreendeudamiento y externalidades financieras; 9.4 Crisis de balanza de pagos & 10 \\
26 & 9 & 9.5 Deuda externa y deuda soberana; 9.6 Dolarización y crisis financieras & 10 \\
\bottomrule
\end{tabularx}
}

%==============================================================================
\parte{Parte I --- Comercio internacional}

%------------------------------------------------------------------------------
\bloque{Bloque 1 --- Comercio mundial y ventaja comparativa ricardiana --- 2 semanas}
{Los hechos del comercio mundial y el primer modelo de la ventaja
comparativa. Cubre los temas 1 y 2 de ECO208, el tema 1 de 1ECO16 y la primera
parte del Módulo 1 de 1ECO83.}

\tema{1.1 El comercio mundial: una visión general}
\lecturas{\KO{1--2}}
\nota{Quién comercia con quién, el tamaño de las economías y la distancia: la
versión intuitiva de la ecuación de gravedad que el tema 3.3 formaliza.}

\tema{1.2 El modelo ricardiano}
\lecturas{\KO{3} \quad \FE{1} (modelo ricardiano)}
\nota{Productividad del trabajo, ventaja comparativa, salarios relativos y
ganancias del comercio.}

\tema{1.3 El modelo ricardiano con un continuo de bienes y su evidencia}
\articulos{Dornbusch, Fischer y Samuelson (1977) \quad Bernhofen y Brown (2004)}
\consulta{Costinot (2009) \quad Bernhofen y Brown (2005)}
\nota{El modelo de Dornbusch--Fischer--Samuelson, base de los modelos
cuantitativos modernos; Japón tras su apertura en el siglo XIX como prueba
directa de la ventaja comparativa.}

%------------------------------------------------------------------------------
\bloque{Bloque 2 --- Dotaciones de factores y distribución del ingreso --- 3 semanas}
{Los modelos de factores específicos y de Heckscher--Ohlin, sus teoremas y su
contrastación. Cubre los temas 2 a 5 de 1ECO16, los temas 3 a 5 de ECO208 y
la sesión 1.2 del Módulo 1 de 1ECO83.}

\tema{2.1 El modelo de factores específicos}
\lecturas{\KO{4} \quad \FE{3} (factores específicos)}
\articulos{Neary (1978)}
\nota{Quién gana y quién pierde con el comercio en el corto plazo, cuando
algunos factores no pueden cambiar de sector.}

\tema{2.2 El modelo de Heckscher--Ohlin y sus teoremas}
\lecturas{\KO{5} \quad \FE{1} (modelo de dos sectores) \quad \FE{2}}
\articulos{Stolper y Samuelson (1941) \quad Rybczynski (1955)}
\nota{Los teoremas de Stolper--Samuelson, de Rybczynski, de Heckscher--Ohlin y
de igualación del precio de los factores, con la dualidad de costos que
Feenstra usa en todo el libro.}

\tema{2.3 Muchos bienes y factores: el modelo de Heckscher--Ohlin--Vanek}
\lecturas{\FE{3}}
\articulos{Dornbusch, Fischer y Samuelson (1980)}
\nota{El contenido factorial del comercio y las generalizaciones de los
teoremas a dimensiones mayores.}

\tema{2.4 La evidencia empírica}
\lecturas{\FE{2} (secciones empíricas)}
\articulos{Leamer (1980) \quad Trefler (1993) \quad Trefler (1995) \quad Davis y Weinstein (2001)}
\consulta{Bowen, Leamer y Sveikauskas (1987) \quad Bernhofen y Brown (2016)}
\nota{La paradoja de Leontief, el comercio «faltante» y las pruebas con
diferencias tecnológicas entre países.}

\tema{2.5 El modelo estándar, los términos de intercambio y los auges de materias primas}
\lecturas{\KO{6}}
\articulos{Corden y Neary (1982)}
\consulta{Corden (1984) \quad Deaton y Laroque (1992) \quad Hausmann y Klinger (2008)}
\nota{Crecimiento sesgado, transferencias y términos de intercambio; la
enfermedad holandesa en una economía pequeña con recursos naturales, el tema
de la sesión 1.3 de 1ECO83 y del tema 8 de ECO208.}

%------------------------------------------------------------------------------
\bloque{Bloque 3 --- Rendimientos crecientes, empresas y gravedad --- 3 semanas}
{La «nueva teoría» del comercio (Krugman) y la «nueva nueva teoría» de las
empresas heterogéneas (Melitz). Cubre los temas 6 y 7 de 1ECO16, el tema 6 de
ECO208 y el Módulo 2 de 1ECO83.}

\tema{3.1 Economías externas de escala y localización internacional}
\lecturas{\KO{7}}
\nota{Concentración geográfica, el efecto del mercado interno y el papel de la
historia en el patrón de especialización.}

\tema{3.2 Competencia monopolística y comercio intraindustrial}
\lecturas{\FE{5} (modelo de Krugman) \quad \HK{7.1--7.4}}
\articulos{Krugman (1979) \quad Krugman (1980)}
\consultalib{Dagenais y Muet (1992), cap.~1}
\nota{Preferencias por la variedad de Dixit--Stiglitz (sílabo de
microeconomía), comercio entre países parecidos y el efecto del mercado
interno.}

\tema{3.3 La ecuación de gravedad}
\lecturas{\FE{5} (ecuación de gravedad)}
\articulos{Anderson y van Wincoop (2003) \quad McCallum (1995)}
\consulta{Anderson (1979) \quad Helpman (1987) \quad Hummels y Levinsohn (1995)}
\nota{Fundamentos teóricos, la resistencia multilateral y el efecto frontera
entre Canadá y Estados Unidos.}

\tema{3.4 Empresas heterogéneas: el modelo de Melitz}
\lecturas{\KO{8}}
\articulos{Melitz (2003) \quad Bernard y Jensen (1999) \quad Bustos (2011)}
\consulta{Khandelwal (2010) \quad Kugler y Verhoogen (2012)}
\nota{Selección de las empresas más productivas en la exportación,
reasignación dentro de la industria y ganancias de productividad. Bustos
estudia la mejora tecnológica de las empresas argentinas con el MERCOSUR.}

\tema{3.5 Comercio de insumos intermedios, deslocalización y salarios}
\lecturas{\FE{4}}
\nota{La fragmentación de la producción y su efecto sobre la prima salarial
por calificación; enlaza con el tema 4.6.}

\tema{3.6 Empresas multinacionales y organización de la empresa}
\lecturas{\FE{11}}
\apoyo{\FE{Apéndice B} (modelos de elección discreta)}
\nota{Inversión extranjera directa horizontal y vertical, contratos
incompletos e integración vertical frente a subcontratación.}

%------------------------------------------------------------------------------
\bloque{Bloque 4 --- Política comercial --- 3 semanas}
{Instrumentos, competencia imperfecta, economía política, acuerdos regionales y
efectos distributivos. Cubre los temas 7 a 9 de ECO208, el tema 8 de 1ECO16,
los Módulos 3 y 4 de 1ECO83 y la primera parte de Helpman y Krugman.}

\tema{4.1 Instrumentos de política comercial y protección efectiva}
\lecturas{\KO{9} \quad \FE{7} \quad \FE{8}}
\apoyo{\HK{2.1--2.6}}
\articulos{Corden (1971)}
\nota{Aranceles, cuotas, subsidios a la exportación y restricciones voluntarias
en equilibrio parcial y general; el arancel óptimo, el \emph{dumping} y la
protección efectiva del Módulo 4 de 1ECO83.}

\tema{4.2 Política comercial con competencia imperfecta}
\lecturas{\HK{3--6}}
\consultalib{McLaren (2013), cap.~4}
\nota{Aranceles frente a un monopolista extranjero, oligopolio de Cournot y
de Bertrand, traslado de beneficios y política comercial estratégica
(Brander--Spencer).}

\tema{4.3 Economía política de la política comercial}
\lecturas{\KO{10} \quad \FE{9}}
\nota{Grupos de interés, la protección en venta de Grossman--Helpman y la
negociación internacional en la OMC.}

\tema{4.4 Acuerdos comerciales regionales}
\lecturas{\FE{6}}
\articulos{Panagariya (2000) \quad Limão (2006) \quad Estevadeordal, Freund y Ornelas (2008) \quad Tovar (2012)}
\consulta{Bown y Tovar (2016) \quad Tovar (2019) \quad Banco Mundial (2020)}
\nota{Creación y desviación de comercio y la discusión sobre si los acuerdos
preferenciales son escalones o tropiezos hacia la liberalización multilateral,
con evidencia latinoamericana.}

\tema{4.5 Comercio, desarrollo y crecimiento}
\lecturas{\KO{11} \quad \FE{10}}
\consultalib{Acemoglu (2009), secc.~19.3--19.7}
\nota{Sustitución de importaciones, promoción de exportaciones, industria
naciente y los modelos de comercio con crecimiento endógeno.}

\tema{4.6 Comercio y desigualdad: evidencia y el caso peruano}
\lecturas{\KO{12} \quad Morales, Pierola y Sánchez-Navarro (2021)}
\articulos{Pavcnik (2017) \quad Autor, Dorn y Hanson (2013)}
\consultalib{Sanborn, Pareja y Quispe (2024)}
\consulta{FMI (2023), \emph{Geo-economic Fragmentation and the Future of Multilateralism}}
\nota{El «\emph{shock} de China» en los mercados laborales locales y el efecto
de la competencia importada sobre la informalidad y los salarios en la
manufactura peruana. Sanborn et al.\ documentan la presencia económica china
en el Perú, de Marcona a Chancay.}

%==============================================================================
\parte{Parte II --- Macroeconomía internacional}

%------------------------------------------------------------------------------
\bloque{Bloque 5 --- La cuenta corriente --- 4 semanas}
{Las cuentas externas y el enfoque intertemporal de la cuenta corriente, de
los modelos de dos períodos de Schmitt-Grohé et al.\ a los de horizonte
infinito de Uribe y Schmitt-Grohé. Retoma el Bloque 3 del sílabo de
macroeconomía con más profundidad. Cubre la sesión 1.3 de 1ECO93 y la
sesión 4.2 de su Módulo 4.}

\tema{5.1 Balanza de pagos y posición de inversión internacional}
\lecturas{\SW{1.1--1.8} \quad \KO{13}}
\nota{Cuenta corriente y financiera, efectos de valoración y la «materia
oscura» de Estados Unidos.}

\tema{5.2 Sostenibilidad de la cuenta corriente}
\lecturas{\SW{2.1--2.4}}
\nota{¿Puede un país tener déficits perpetuos de balanza comercial o de cuenta
corriente? Restricción presupuestaria intertemporal y condición de no juego
de Ponzi.}

\tema{5.3 La cuenta corriente en una economía de dotación}
\lecturas{\SW{3.1--3.7} \quad \US{2.1--2.2}}
\apoyo{\VE{1}}
\nota{Choques temporales y permanentes de producción, de términos de
intercambio y de tasa de interés. El modelo de horizonte infinito de \US{2}
muestra que la cuenta corriente es contracíclica cuando los choques son
persistentes.}

\tema{5.4 Inversión y cuenta corriente}
\lecturas{\SW{4.1--4.4} \quad \US{3.1--3.5}}
\nota{Decisión de inversión, curva de ahorro, choques de productividad
permanentes y temporales y costos de ajuste del capital.}

\tema{5.5 Incertidumbre y prueba del enfoque intertemporal}
\lecturas{\SW{5.1--5.2} \quad \US{2.3--2.5}}
\nota{Ahorro precautorio, la gran moderación y las pruebas econométricas del
enfoque intertemporal de la cuenta corriente.}

\tema{5.6 Ajuste externo en economías pequeñas y grandes}
\lecturas{\SW{6.1--6.6}}
\nota{Primas de riesgo país, la hipótesis del exceso de ahorro mundial y los
controles de capital óptimos en un modelo de dos países.}

\tema{5.7 Déficits gemelos y equivalencia ricardiana}
\lecturas{\SW{7.1--7.6}}
\nota{El sector público en la economía abierta y las condiciones bajo las que
el déficit fiscal se traslada a la cuenta corriente.}

\tema{5.8 Imperfecciones del mercado de capitales}
\lecturas{\VE{2}}
\consulta{Obstfeld y Rogoff (1996), caps.~1--2}
\nota{Restricciones de endeudamiento, primas de riesgo crecientes y su efecto
sobre el consumo y la cuenta corriente en economías emergentes.}

%------------------------------------------------------------------------------
\bloque{Bloque 6 --- Integración financiera y tipo de cambio real --- 2 semanas}
{Cuán integrados están los mercados de capital y qué determina el tipo de
cambio real. Retoma los temas 3.3 y 3.4 del sílabo de macroeconomía y el tema
8.1 del de economía financiera.}

\tema{6.1 Movilidad del capital y paridades de tasas de interés}
\lecturas{\SW{8.1--8.3}}
\articulos{Feldstein y Horioka (1980)}
\consultalib{\KO{14--15} (leídos en el sílabo de economía financiera)}
\nota{Correlaciones entre ahorro e inversión, diferenciales de tasas de interés
y el enigma de la paridad descubierta.}

\tema{6.2 Determinantes del tipo de cambio real}
\lecturas{\SW{9.1--9.6}}
\articulos{Taylor y Taylor (2004) \quad Gopinath y Rigobon (2008)}
\consultalib{\KO{16} (leído en el sílabo de economía financiera)}
\nota{Ley de un solo precio, paridad del poder de compra, bienes no
transables, el efecto Balassa--Samuelson, barreras comerciales y
microfundamentos de los índices de precios. Es el tema 9 de 1ECO16.}

\tema{6.3 Bienes no transables y precios intertemporales}
\lecturas{\US{7.9--7.10}}
\apoyo{\VE{3--4}}
\nota{El tipo de cambio real como precio relativo de los no transables y su
papel como precio intertemporal en una economía de dotación.}

\tema{6.4 El modelo TNT y las paradas súbitas}
\lecturas{\SW{10.1--10.5} \quad \US{7.11--7.12}}
\nota{Choques de demanda agregada y ajuste del tipo de cambio real ante las
paradas súbitas de Argentina (2001), Chile (años ochenta) y Turquía (2001).}

%------------------------------------------------------------------------------
\bloque{Bloque 7 --- Ciclos económicos en economías abiertas y emergentes --- 3 semanas}
{Los capítulos 1 y 4 a 7 de Uribe y Schmitt-Grohé. Cubre los temas de economía
abierta de 1ECO74 (ciclos reales en economía abierta y términos de
intercambio), las sesiones 1.1 y 1.2 de 2ECO03 y la sección 4 de ECO338.}

\tema{7.1 Hechos estilizados del ciclo en el mundo}
\lecturas{\US{1.1--1.8}}
\nota{Diez regularidades: la volatilidad del consumo y la contraciclicidad de
la balanza comercial, en países pobres, emergentes y ricos.}

\tema{7.2 El modelo RBC de economía pequeña y abierta}
\lecturas{\US{4.1--4.10}}
\articulos{Mendoza (1991) \quad Schmitt-Grohé y Uribe (2003)}
\nota{Descentralización, calibración, aproximación de la dinámica, mercados
completos frente a incompletos y las formas de inducir estacionariedad (el
problema de «cerrar» los modelos de economía pequeña). Las sesiones 1.1 y 1.2
de 2ECO03 comparan los dos tipos de mercados.}

\tema{7.3 Ciclos en economías emergentes}
\lecturas{\US{5.1--5.7}}
\articulos{Aguiar y Gopinath (2007) \quad García-Cicco, Pancrazi y Uribe (2010)}
\nota{¿El ciclo es la tendencia? Choques permanentes de productividad frente a
fricciones financieras, con estimación bayesiana sobre un siglo de datos.}

\tema{7.4 Choques de tasa de interés}
\lecturas{\US{6.1--6.7}}
\nota{Un modelo empírico y uno teórico de los choques a la tasa de interés
internacional y al riesgo país, y los factores de riesgo global.}

\tema{7.5 Términos de intercambio}
\lecturas{\US{7.1--7.8}}
\articulos{Schmitt-Grohé y Uribe (2018)}
\consulta{Castillo y Rojas (2014) \quad Rodríguez, Villanueva y Castillo (2018)}
\nota{El efecto Harberger--Laursen--Metzler y la importancia de los choques de
términos de intercambio en un modelo RBC con transables y no transables;
evidencia para el Perú y América Latina.}
\nota{Para un importador neto de energía, un alza del precio del petróleo opera
como un deterioro de los términos de intercambio. La versión de economía
cerrada, con la energía como insumo de la producción, es el RBC con energía del
tema 12.6 del sílabo de macroeconomía.}

\tema{7.6 Ciclos internacionales y los enigmas de la macroeconomía abierta}
\articulos{Backus, Kehoe y Kydland (1992) \quad Backus y Smith (1993) \quad Obstfeld y Rogoff (2000)}
\consulta{Meese y Rogoff (1983) \quad Itskhoki y Mukhin (2021) \quad Gabaix y Maggiori (2015)}
\nota{Comovimiento internacional, el enigma de Backus--Smith, Feldstein--Horioka
y la desconexión del tipo de cambio: la sección 4 de ECO338.}

%------------------------------------------------------------------------------
\bloque{Bloque 8 --- Dinero, tipo de cambio y política cambiaria --- 3 semanas}
{La parte nominal: tipo de cambio, rigideces, regímenes cambiarios y la
estabilización de la inflación. Cubre el Módulo 2 de 1ECO93, el Módulo 2 de
2ECO03 y los temas 10 y 11 de 1ECO16.}

\tema{8.1 Dinero, inflación y tipo de cambio nominal}
\lecturas{\SW{13.1--13.5}}
\apoyo{\VE{5--6}}
\nota{Teoría cuantitativa, déficit fiscal e inflación, tipo de cambio fijo y
crecimiento constante del dinero, impuesto inflacionario y el enfoque
monetario de la balanza de pagos.}

\tema{8.2 Producción a corto plazo y áreas monetarias óptimas}
\lecturas{\KO{17} \quad \KO{21}}
\articulos{Bayoumi (1994) \quad Rose (2000)}
\consulta{Barro y Tenreyro (2007) \quad Mongelli (2002)}
\nota{El modelo DD--AA de producción y tipo de cambio a corto plazo y la teoría
de las áreas monetarias óptimas aplicada al euro.}

\tema{8.3 Rigidez salarial, desempleo y política cambiaria óptima}
\lecturas{\SW{11.1--11.2} \quad \US{8.1--8.4} \quad \US{8.8--8.10}}
\apoyo{\US{8.11--8.15}}
\nota{Rigidez a la baja de los salarios nominales, el costo de las paridades
fijas y la política cambiaria óptima, con la periferia europea en 2008--2011
como caso.}

\tema{8.4 El modelo neokeynesiano de economía abierta}
\lecturas{\GA{7}}
\articulos{De Paoli (2009)}
\consultalib{Walsh (2010), secc.~9.1--9.5}
\nota{La economía pequeña y abierta de Galí y Monacelli: curva de Phillips,
IS dinámica y reglas de política. El sílabo de macroeconomía lo cita solo
como consulta (tema 13.1); Walsh añade el modelo de dos países de
Obstfeld--Rogoff y la coordinación de políticas.}

\tema{8.5 Regímenes cambiarios y miedo a flotar}
\lecturas{\VE{11}}
\articulos{Calvo y Reinhart (2002) \quad Fischer (2001) \quad Calvo y Mishkin (2003) \quad Obstfeld y Rogoff (1995)}
\consultalib{\KO{18--19} (leídos en el sílabo de economía financiera)}
\nota{Elección óptima del régimen cambiario, la visión bipolar y por qué las
economías emergentes temen flotar.}

\tema{8.6 Estabilización de la inflación alta}
\lecturas{\VE{12--13}}
\nota{Anclas nominales y reales, programas basados en el tipo de cambio y en
el dinero, y el ciclo de auge y recesión de las estabilizaciones; marco para
la hiperinflación peruana de 1988--1990.}

%------------------------------------------------------------------------------
\bloque{Bloque 9 --- Flujos de capital, crisis y deuda soberana --- 3 semanas}
{Los capítulos 9 a 11 de Uribe y Schmitt-Grohé y la última parte de Végh.
Cubre los Módulos 1, 3 y 4 de 1ECO93, los Módulos 1 y 3 de 2ECO03, el tema
12 de 1ECO16 y la dolarización de 1ECO74.}

\tema{9.1 Flujos de capital y paradas súbitas}
\lecturas{\VE{14}}
\articulos{Cavallo, Powell, Pedemonte y Tavella (2015) \quad Forbes y Warnock (2021)}
\consultalib{Acemoglu (2009), secc.~19.1--19.2}
\consulta{Koepke (2019) \quad Ahmed y Zlate (2014)}
\nota{Determinantes de los flujos hacia economías emergentes, olas de entradas
y salidas y las respuestas de política.}

\tema{9.2 Controles de capital}
\lecturas{\US{9.1--9.9}}
\consulta{Ostry, Ghosh, Chamon y Qureshi (2012)}
\nota{Política fiscal de primer óptimo con tipo de cambio fijo, controles de
capital como segundo óptimo y su carácter prudencial.}

\tema{9.3 Sobreendeudamiento y externalidades financieras}
\lecturas{\US{10.1--10.2}}
\articulos{Mendoza (2010) \quad Bianchi (2011)}
\consulta{Bianchi y Mendoza (2020) \quad Aoki, Benigno y Kiyotaki (2021)}
\nota{Credibilidad imperfecta, restricciones de colateral y la externalidad
pecuniaria que justifica la política macroprudencial (sesiones 1.3--1.6 de
2ECO03 y 3.2 de 1ECO93).}

\tema{9.4 Crisis de balanza de pagos}
\lecturas{\SW{13.6} \quad \VE{16}}
\articulos{Agénor, Bhandari y Flood (1992)}
\consultalib{\KO{20} y \KO{22} (leídos en el sílabo de economía financiera)}
\consulta{Obstfeld, \emph{The Logic of Currency Crises}}
\nota{Crisis de primera generación (ataques especulativos, Perú 1985--1990),
de segunda generación (la libra esterlina en 1992) y de tercera (Asia 1997).}

\tema{9.5 Deuda externa y deuda soberana}
\lecturas{\SW{12.1--12.4} \quad \US{11.1--11.6}}
\apoyo{\US{11.7--11.9}}
\nota{La crisis de la deuda de los ochenta, esquemas de reducción de deuda, el
costo del impago y el modelo de Eaton--Gersovitz.}

\tema{9.6 Dolarización y crisis financieras}
\lecturas{\VE{15} \quad \VE{17}}
\articulos{Christiano, Dalgic y Nurbekyan (2021) \quad Carrasco y Florián (2025)}
\consulta{Carrasco, Florián y Nivín (2019) \quad Céspedes, Chang y Velasco (2017)}
\consultalib{Castillo, Montoro y Tuesta (2009) \quad Montoro y Ortiz (2016) \quad Gopinath (2019)}
\nota{Determinantes de la dolarización financiera, crisis gemelas, intervención
cambiaria y acumulación de reservas en una economía semidolarizada como la
peruana (Módulo 3 de 2ECO03).}
\nota{Castillo, Montoro y Tuesta (2009) dan el modelo de equilibrio general con
dolarización para el Perú; Montoro y Ortiz (2016), un modelo de la intervención
cambiaria y su interacción con la regla de tasas; Gopinath (2019), el marco de
políticas integradas del FMI. Los tres los usa el curso de Política Monetaria
del BCRP.}

%------------------------------------------------------------------------------
\newpage
\section*{Anexo A. De los capítulos a los bloques}
\addcontentsline{toc}{section}{Anexo A. De los capítulos a los bloques}

Capítulos de cada libro que entran en el plan y temas donde se leen. Lo que
queda fuera se detalla en el Anexo B.

\newcommand{\mapa}[2]{%
\subsection*{#1}
{\small\filasalternas
\noindent\begin{tabularx}{\textwidth}{>{\raggedright\arraybackslash}p{1.1cm}X>{\raggedright\arraybackslash}p{3.4cm}}
\toprule
\rowcolor{tonoOscuro}\cab{Cap.} & \cab{Título} & \cab{Temas} \\
\midrule
#2
\bottomrule
\end{tabularx}
}}

\mapa{A.1 Feenstra, \emph{Advanced International Trade}}{%
1 & Preliminaries: Two-Sector Models & 1.2, 2.2 \\
2 & The Heckscher-Ohlin Model & 2.2, 2.4 \\
3 & Many Goods and Factors & 2.1, 2.3 \\
4 & Trade in Intermediate Inputs and Wages & 3.5 \\
5 & Increasing Returns and the Gravity Equation & 3.2, 3.3 \\
6 & Gains from Trade and Regional Agreements & 4.4 \\
7 & Import Tariffs and Dumping & 4.1 \\
8 & Import Quotas and Export Subsidies & 4.1 \\
9 & Political Economy of Trade Policy & 4.3 \\
10 & Trade and Endogenous Growth & 4.5 \\
11 & Multinationals and Organization of the Firm & 3.6 \\
B & Discrete Choice Models & 3.6 \\
}

\mapa{A.2 Krugman, Obstfeld y Melitz, \emph{Economía internacional}}{%
1--2 & Introducción; Comercio mundial: una visión general & 1.1 \\
3 & Productividad del trabajo y ventaja comparativa: el modelo ricardiano & 1.2 \\
4 & Factores específicos y distribución de la renta & 2.1 \\
5 & Recursos y comercio: el modelo de Heckscher-Ohlin & 2.2 \\
6 & El modelo estándar de comercio & 2.5 \\
7 & Economías externas de escala y localización internacional de la producción & 3.1 \\
8 & Las empresas en la economía global & 3.4 \\
9 & Los instrumentos de la política comercial & 4.1 \\
10 & La economía política de la política comercial & 4.3 \\
11 & La política comercial en los países en desarrollo & 4.5 \\
12 & Controversias de la política comercial & 4.6 \\
13 & La contabilidad nacional y la balanza de pagos & 5.1 \\
17 & La producción nacional y el tipo de cambio a corto plazo & 8.2 \\
21 & Áreas monetarias óptimas y el euro & 8.2 \\
}

\mapa{A.3 Helpman y Krugman, \emph{Trade Policy and Market Structure}}{%
2 & Trade Policy under Perfect Competition (2.1--2.6) & 4.1 \\
3--4 & Protection and Domestic Market Power; Foreign Market Power & 4.2 \\
5--6 & Strategic Export Policy; Strategic Import Policy & 4.2 \\
7 & Intraindustry Trade (7.1--7.4) & 3.2 \\
}

\mapa{A.4 Schmitt-Grohé, Uribe y Woodford, \emph{International Macroeconomics}}{%
1 & Global Imbalances & 5.1 \\
2 & Current Account Sustainability & 5.2 \\
3 & A Theory of Current Account Determination & 5.3 \\
4 & Current Account Determination in a Production Economy & 5.4 \\
5 & Uncertainty and the Current Account & 5.5 \\
6 & External Adjustment in Small and Large Economies & 5.6 \\
7 & Twin Deficits & 5.7 \\
8 & International Capital Market Integration & 6.1 \\
9 & Determinants of the Real Exchange Rate & 6.2 \\
10 & Aggregate Demand Shocks and Real Exchange Rates & 6.4 \\
11 & Exchange Rate Policy and Unemployment & 8.3 \\
12 & The Macroeconomics of External Debt & 9.5 \\
13 & Monetary Policy and Nominal Exchange Rate Determination & 8.1, 9.4 \\
}

\mapa{A.5 Uribe y Schmitt-Grohé, \emph{Open Economy Macroeconomics}}{%
1 & Business-Cycle Facts Around the World (1.1--1.8) & 7.1 \\
2 & An Open Endowment Economy & 5.3, 5.5 \\
3 & An Open Economy with Capital & 5.4 \\
4 & The Small-Open-Economy Real-Business-Cycle Model (4.1--4.10) & 7.2 \\
5 & Emerging-Country Business Cycles Through the Lens of the SOE-RBC Model & 7.3 \\
6 & Interest-Rate Shocks & 7.4 \\
7 & Tradable Goods, Nontradable Goods, the Terms of Trade, and the Real Exchange Rate & 6.3, 6.4, 7.5 \\
8 & Nominal Rigidity, Exchange Rates, and Unemployment & 8.3 \\
9 & Exchange Rate Policy and Capital Controls & 9.2 \\
10 & Overborrowing & 9.3 \\
11 & Sovereign Debt & 9.5 \\
}

\mapa{A.6 Végh, \emph{Open Economy Macroeconomics in Developing Countries}}{%
1--2 & The Basic Intertemporal Model; Capital Market Imperfections & 5.3, 5.8 \\
3--4 & Intertemporal Prices; Nontradable Goods and Relative Prices & 6.3 \\
5--6 & The Basic Monetary Model; The Monetary Approach to the Balance of Payments & 8.1 \\
11 & Optimal Exchange Rate Regimes & 8.5 \\
12--13 & Real Anchors; Stopping High Inflation & 8.6 \\
14 & Capital Inflows & 9.1 \\
15 & Dollarization & 9.6 \\
16 & Balance-of-Payments Crises & 9.4 \\
17 & Financial Crises & 9.6 \\
}

\mapa{A.7 Galí, \emph{Monetary Policy, Inflation, and the Business Cycle}}{%
7 & Monetary Policy and the Open Economy & 8.4 \\
}

%------------------------------------------------------------------------------
\section*{Anexo B. Lo que no entra en el plan}
\addcontentsline{toc}{section}{Anexo B. Lo que no entra en el plan}

{\small\filasalternas
\noindent\begin{tabularx}{\textwidth}{>{\raggedright\arraybackslash}p{4.6cm}X}
\toprule
\rowcolor{tonoOscuro}\cab{Materia} & \cab{Capítulos y secciones} \\
\midrule
Finanzas internacionales ya cubiertas & \KO{14--16}, \KO{18--20} y \KO{22}: Bloque 8 del sílabo de economía financiera; aquí, solo consulta \\
Métodos numéricos & \US{4.11--4.16} (aproximaciones de primer orden, momentos, impulso-respuesta y código de Matlab); \HK{8} (cuantificación); Piermartini y Teh (2005) \\
Índices de precios y productividad & \FE{Apéndice A} \\
Política macroeconómica de corto plazo en Végh & \VE{7--10} (política temporal, precios rígidos, política de tasas de interés y política fiscal y monetaria óptima) (Anexo D.2) \\
Comercio y política con enfoque de casos & McLaren (2013): solo el cap.~4 como consulta \\
Temas de coyuntura & Los informes y notas de organismos del Módulo 3 de 1ECO83 (BID, CAF, CEPAL, FMI, OMC, MINCETUR), salvo FMI (2023) como consulta \\
\bottomrule
\end{tabularx}
}

\medskip
\noindent\destacado{Archivos de la carpeta que no se usaron}: Obstfeld y Rogoff,
\emph{Foundations of International Macroeconomics} (carpeta de
macroeconomía), por ser una imagen sin texto; se cita como consulta porque lo
usa 1ECO93. La copia de Krugman et al.\ terminada en
\texttt{.crdownload} es una descarga incompleta; la 9.ª edición y la versión de
2006 quedan sustituidas por la 10.ª. Tampoco se usaron Dagenais y Muet (eds.),
\emph{International Trade Modelling} (salvo su cap.~1 como
consulta), por tratarse de una colección de estudios empíricos de comienzos de los años
noventa; Shaikh (2009), \emph{Teorías del comercio internacional}, cuyo texto
escaneado es de mala calidad (Anexo D.4); y el trabajo de Liu (2023), un
informe de estudiante.

\medskip
\noindent\destacado{Versiones de Uribe y Schmitt-Grohé}: la carpeta de economía
internacional tiene un borrador de agosto de 2014 (11 capítulos, 567 páginas)
con un nombre de archivo que indica la edición de 2017; el plan usa el
borrador de marzo de 2015 de la carpeta de macroeconomía, más completo
(capítulo 7 ampliado y secciones 11.7--11.9) y el mismo que cita el sílabo de
macroeconomía.

%------------------------------------------------------------------------------
\section*{Anexo C. Programas utilizados}
\addcontentsline{toc}{section}{Anexo C. Programas utilizados}

{\small\filasalternas
\noindent\begin{tabularx}{\textwidth}{>{\raggedright\arraybackslash}p{5.6cm}X}
\toprule
\rowcolor{tonoOscuro}\cab{Programa} & \cab{Temas del plan} \\
\midrule
ECO208, Economía Internacional I (pregrado, 2025-2 y 2026-1, J.~D.~Gallardo) & Temas 1 a 9 del programa: Bloques 1 a 4; tema 10 (balanza de pagos y tipo de cambio), 5.1 y 6.1 \\
1ECO16, Economía Internacional 2 (pregrado, 2025-2, P.~Tovar) & Temas 1 a 7: Bloques 1 a 3; tema 8, 4.4; temas 9 a 11, 6.2, 8.2 y 8.5; tema 12, 9.4 \\
1ECO83, Economía Internacional (Maestría en Economía, 2025-2, R.~Castellares, A.~Fairlie y D.~Schydlowsky) & Módulo 1, Bloques 1 y 2; módulo 2, Bloque 3; módulo 3, 4.4 y 4.6; módulo 4, 4.1 \\
1ECO93, Finanzas Internacionales (Maestría en Economía, 2025-2) & Módulo 1, 5.2 y 9.4; módulo 2, 8.2; módulo 3, 9.1 y 9.3; módulo 4, 5.1, 5.8, 9.5 y 9.6 \\
2ECO03, Macroeconomía Avanzada (Maestría en Economía, 2025-2) & Módulo 1, 7.2, 9.2 y 9.3; módulo 2, 8.4; módulo 3, 9.6 \\
1ECO74, Macroeconomía Intermedia (Maestría en Economía, 2026-1) & Ciclos reales en economía abierta y términos de intercambio, Bloque 7; dolarización, 9.6 \\
ECO338, Tópicos de Macroeconomía Avanzada~2 (pregrado, 2026-2, C.~Vásquez) & Sección 4 (macroeconomía internacional), 7.6 \\
\bottomrule
\end{tabularx}
}

\medskip
\noindent\destacado{Programas de la carpeta que no se usaron}: 1FIN27 (Finanzas
Internacionales, 2026-1), porque la versión 2025-2 ya está integrada en el
sílabo de economía financiera. El módulo 4 de 2ECO03 (búsqueda y
emparejamiento) y las unidades de economía cerrada de 1ECO74 ya están en el
sílabo de macroeconomía.

%------------------------------------------------------------------------------
\section*{Anexo D. Material opcional}
\addcontentsline{toc}{section}{Anexo D. Material opcional}

Todo lo de este anexo es opcional: no entra en el cómputo de horas ni modifica
el cronograma.

\subsection*{D.1 Comercio y política comercial con casos}
McLaren (2013), \emph{International Trade}: cada capítulo parte de un problema
de política real (autosuficiencia alimentaria, el TLCAN, el medio ambiente,
el trabajo infantil, el déficit comercial y el renminbi).

\subsection*{D.2 Política macroeconómica de corto plazo en economías en desarrollo}
\VE{7--10}: política temporal, precios rígidos, defensa de la moneda con la
tasa de interés y política fiscal y monetaria óptima en la economía abierta.
Momento sugerido: tras el Bloque 8.

\subsection*{D.3 Modelos cuantitativos de política comercial}
Piermartini y Teh (2005), \emph{Demystifying Modelling Methods for Trade
Policy} (OMC); \HK{8}. Requiere métodos numéricos que corresponden a otro
curso.

\subsection*{D.4 Visiones heterodoxas del comercio}
Shaikh (2009), \emph{Teorías del comercio internacional}: crítica de la ventaja
comparativa desde la teoría de la ventaja absoluta y los costos reales.

\subsection*{D.5 Coyuntura y fragmentación geoeconómica}
Los informes del Módulo 3 de 1ECO83: Banco Mundial (2020), FMI (2023), OMC
(2024) y los documentos sobre los acuerdos comerciales del Perú y la Alianza
del Pacífico.

%------------------------------------------------------------------------------
\section*{Anexo E. Referencias}
\addcontentsline{toc}{section}{Anexo E. Referencias}

\subsection*{E.1 Textos del plan}
{\small
\begin{itemize}[leftmargin=1.4em,itemsep=1pt,label={}]
\item Feenstra, R. C. (2004). \emph{Advanced International Trade: Theory and Evidence}. Princeton University Press.
\item Galí, J. (2008). \emph{Monetary Policy, Inflation, and the Business Cycle: An Introduction to the New Keynesian Framework}. Princeton University Press (carpeta de macroeconomía).
\item Helpman, E. y P. R. Krugman (1989). \emph{Trade Policy and Market Structure}. The MIT Press.
\item Krugman, P. R., M. Obstfeld y M. J. Melitz (2016). \emph{Economía internacional: teoría y política}, 10.ª ed. Pearson.
\item Morales, F., M. D. Pierola y D. Sánchez-Navarro (2021). Import Competition in the Manufacturing Sector in Peru: Its Impact on Informality and Wages. \emph{Economía}, 44(88), 45--75.
\item Schmitt-Grohé, S., M. Uribe y M. Woodford (2016). \emph{International Macroeconomics}. Borrador del 4 de mayo de 2016; publicado como \emph{International Macroeconomics: A Modern Approach}, Princeton University Press, 2022.
\item Uribe, M. y S. Schmitt-Grohé (2015). \emph{Open Economy Macroeconomics}. Borrador del 17 de marzo de 2015; publicado por Princeton University Press, 2017 (carpeta de macroeconomía).
\item Végh, C. A. (2013). \emph{Open Economy Macroeconomics in Developing Countries}. The MIT Press.
\end{itemize}
}

\subsection*{E.2 Libros de consulta de la carpeta (sin horas)}
{\small
\begin{itemize}[leftmargin=1.4em,itemsep=1pt,label={}]
\item Acemoglu, D. (2009). \emph{Introduction to Modern Economic Growth}. Princeton University Press (carpeta de desarrollo y crecimiento).
\item Dagenais, M. G. y P.-A. Muet (eds.) (1992). \emph{International Trade Modelling}. Chapman \& Hall.
\item McLaren, J. (2013). \emph{International Trade: Economic Analysis of Globalization and Policy}. Wiley.
\item Piermartini, R. y R. Teh (2005). \emph{Demystifying Modelling Methods for Trade Policy}. Discussion Paper 10, Organización Mundial del Comercio.
\item Sanborn, C. A., A. Pareja y D. Quispe (2024). \emph{De Marcona a Chancay: la presencia económica y empresarial china en el Perú, 1992--2023}. Documento de Investigación 28, Universidad del Pacífico.
\item Shaikh, A. (2009). \emph{Teorías del comercio internacional}.
\item Walsh, C. E. (2010). \emph{Monetary Theory and Policy}, 3.ª ed. The MIT Press (carpeta de macroeconomía).
\end{itemize}
}

\subsection*{E.3 Referencias tomadas de los programas}
{\small
\begin{itemize}[leftmargin=1.4em,itemsep=1pt,label={}]
\item Agénor, P.-R., J. Bhandari y R. Flood (1992). Speculative Attacks and Models of Balance of Payments Crises. \emph{IMF Staff Papers}, 39(2).
\item Aguiar, M. y G. Gopinath (2007). Emerging Market Business Cycles: The Cycle Is the Trend. \emph{Journal of Political Economy}, 115(1), 69--102.
\item Ahmed, S. y A. Zlate (2014). Capital Flows to Emerging Market Economies: A Brave New World? \emph{Journal of International Money and Finance}, 48, 221--248.
\item Anderson, J. (1979). A Theoretical Foundation for the Gravity Equation. \emph{American Economic Review}, 69(1), 106--116.
\item Anderson, J. y E. van Wincoop (2003). Gravity with Gravitas: A Solution to the Border Puzzle. \emph{American Economic Review}, 93(1), 170--192.
\item Aoki, K., G. Benigno y N. Kiyotaki (2021). Monetary and Financial Policies for Emerging Market Economies. Mimeo.
\item Autor, D., D. Dorn y G. Hanson (2013). The China Syndrome: Local Labor Market Effects of Import Competition in the United States. \emph{American Economic Review}, 103(6), 2121--2168.
\item Backus, D. K., P. J. Kehoe y F. E. Kydland (1992). International Real Business Cycles. \emph{Journal of Political Economy}, 100, 745--775.
\item Backus, D. K. y G. W. Smith (1993). Consumption and Real Exchange Rates in Dynamic Economies with Non-Traded Goods. \emph{Journal of International Economics}, 35, 297--316.
\item Banco Mundial (2020). The Evolution of Deep Trade Agreements. Policy Research Working Paper 9283.
\item Barro, R. y S. Tenreyro (2007). Economic Effects of Currency Unions. \emph{Economic Inquiry}, 45(1), 1--23.
\item Bayoumi, T. (1994). A Formal Model of Optimum Currency Areas. IMF Working Paper 94/42.
\item Bernard, A. y J. B. Jensen (1999). Exceptional Exporter Performance: Cause, Effect, or Both? \emph{Journal of International Economics}, 47(1), 1--25.
\item Bernhofen, D. y J. Brown (2004). A Direct Test of the Theory of Comparative Advantage: The Case of Japan. \emph{Journal of Political Economy}, 112(1), 48--67.
\item Bernhofen, D. y J. Brown (2005). An Empirical Assessment of the Comparative Advantage Gains from Trade: Evidence from Japan. \emph{American Economic Review}, 95(1), 208--225.
\item Bernhofen, D. y J. Brown (2016). Testing the General Validity of the Heckscher-Ohlin Theorem. \emph{American Economic Journal: Microeconomics}, 8(4), 54--90.
\item Bianchi, J. (2011). Overborrowing and Systemic Externalities in the Business Cycle. \emph{American Economic Review}, 101(7), 3400--3426.
\item Bianchi, J. y E. Mendoza (2020). A Fisherian Approach to Financial Crises: Lessons from the Sudden Stops Literature. \emph{Review of Economic Dynamics}, 37, 254--283.
\item Bowen, H., E. Leamer y L. Sveikauskas (1987). Multicountry, Multifactor Tests of the Factor Abundance Theory. \emph{American Economic Review}, 77(5), 791--809.
\item Bown, C. y P. Tovar (2016). Preferential Liberalization, Antidumping, and Safeguards: Stumbling Block Evidence from MERCOSUR. \emph{Economics \& Politics}, 28(3), 262--294.
\item Bustos, P. (2011). Trade Liberalization, Exports, and Technology Upgrading: Evidence on the Impact of MERCOSUR on Argentinian Firms. \emph{American Economic Review}, 101(1), 304--340.
\item Calvo, G. y F. Mishkin (2003). The Mirage of Exchange Rate Regimes for Emerging Market Economies. \emph{Journal of Economic Perspectives}, 17(4), 99--118.
\item Calvo, G. y C. Reinhart (2002). Fear of Floating. \emph{Quarterly Journal of Economics}, 117(2), 379--408.
\item Carrasco, A. y D. Florián (2025). External Shocks and FX Intervention Policy in Financially Dollarized Economies. \emph{Journal of Macroeconomics}, 84, 103672.
\item Carrasco, A., D. Florián y R. Nivín (2019). SFX Interventions, Financial Intermediation, and External Shocks in Emerging Economies. Documento de Trabajo 2019-022, BCRP.
\item Castillo, P. y Y. Rojas (2014). Términos de intercambio y productividad total de factores: evidencia empírica de los mercados emergentes de América Latina. \emph{Revista Estudios Económicos}, 28, 27--46. BCRP.
\item Cavallo, E., A. Powell, M. Pedemonte y P. Tavella (2015). A New Taxonomy of Sudden Stops. \emph{Journal of International Money and Finance}, 51, 47--70.
\item Céspedes, L. F., R. Chang y A. Velasco (2017). Financial Intermediation, Real Exchange Rates, and Unconventional Policies in an Open Economy. \emph{Journal of International Economics}, 108(S1), S76--S86.
\item Christiano, L., H. Dalgic y A. Nurbekyan (2021). Financial Dollarization: Efficient Intranational Risk Sharing or Prescription for Disaster? NBER Working Paper 29034.
\item Corden, W. M. (1971). The Substitution Problem in the Theory of Effective Protection. \emph{Journal of International Economics}, 1(1), 37--57.
\item Corden, W. M. (1984). Booming Sector and Dutch Disease Economics: Survey and Consolidation. \emph{Oxford Economic Papers}, 36(3), 359--380.
\item Corden, W. M. y J. P. Neary (1982). Booming Sector and De-Industrialisation in a Small Open Economy. \emph{Economic Journal}, 92(368), 825--848.
\item Costinot, A. (2009). On the Origins of Comparative Advantage. \emph{Journal of International Economics}, 77, 255--264.
\item Davis, D. y D. Weinstein (2001). An Account of Global Factor Trade. \emph{American Economic Review}, 91(5), 1423--1453.
\item De Paoli, B. (2009). Monetary Policy and Welfare in a Small Open Economy. \emph{Journal of International Economics}, 77(1), 11--22.
\item Deaton, A. y G. Laroque (1992). On the Behaviour of Commodity Prices. \emph{Review of Economic Studies}, 59(1), 1--23.
\item Dornbusch, R., S. Fischer y P. A. Samuelson (1977). Comparative Advantage, Trade and Payments in a Ricardian Model with a Continuum of Goods. \emph{American Economic Review}, 67(5), 823--839.
\item Dornbusch, R., S. Fischer y P. A. Samuelson (1980). Heckscher-Ohlin Trade Theory with a Continuum of Goods. \emph{Quarterly Journal of Economics}, 95(2), 203--224.
\item Estevadeordal, A., C. Freund y E. Ornelas (2008). Does Regionalism Affect Trade Liberalization Toward Nonmembers? \emph{Quarterly Journal of Economics}, 123(4), 1531--1575.
\item Feldstein, M. y C. Horioka (1980). Domestic Saving and International Capital Flows. \emph{Economic Journal}, 90, 314--329.
\item Fischer, S. (2001). Exchange Rate Regimes: Is the Bipolar View Correct? \emph{Journal of Economic Perspectives}, 15(3), 3--24.
\item FMI (2023). Geo-economic Fragmentation and the Future of Multilateralism. Staff Discussion Note 2023/001.
\item Forbes, K. y F. Warnock (2021). Capital Flow Waves---or Ripples? Extreme Capital Flow Movements since the Crisis. \emph{Journal of International Money and Finance}, 116.
\item Gabaix, X. y M. Maggiori (2015). International Liquidity and Exchange Rate Dynamics. \emph{Quarterly Journal of Economics}, 130(3), 1369--1420.
\item García-Cicco, J., R. Pancrazi y M. Uribe (2010). Real Business Cycles in Emerging Countries? \emph{American Economic Review}, 100, 2510--2531.
\item Gopinath, G. y R. Rigobon (2008). Sticky Borders. \emph{Quarterly Journal of Economics}, 123(2), 531--575.
\item Hausmann, R. y B. Klinger (2008). A Growth Diagnosis in Peru. Harvard University, mimeo.
\item Helpman, E. (1987). Imperfect Competition and International Trade: Evidence from Fourteen Industrial Countries. \emph{Journal of the Japanese and International Economies}, 1, 62--81.
\item Hummels, D. y J. Levinsohn (1995). Monopolistic Competition and International Trade: Reconsidering the Evidence. \emph{Quarterly Journal of Economics}, 110(3), 799--836.
\item Itskhoki, O. y D. Mukhin (2021). Exchange Rate Disconnect in General Equilibrium. \emph{Journal of Political Economy}, 129(8), 2183--2232.
\item Khandelwal, A. (2010). The Long and Short (of) Quality Ladders. \emph{Review of Economic Studies}, 77(4), 1450--1476.
\item Koepke, R. (2019). What Drives Capital Flows to Emerging Markets? A Survey of the Empirical Literature. \emph{Journal of Economic Surveys}, 33, 516--540.
\item Krugman, P. (1979). Increasing Returns, Monopolistic Competition, and International Trade. \emph{Journal of International Economics}, 9(4), 469--479.
\item Krugman, P. (1980). Scale Economies, Product Differentiation, and the Pattern of Trade. \emph{American Economic Review}, 70(5), 950--959.
\item Kugler, M. y E. Verhoogen (2012). Prices, Plant Size, and Product Quality. \emph{Review of Economic Studies}, 79(1), 307--339.
\item Leamer, E. (1980). The Leontief Paradox, Reconsidered. \emph{Journal of Political Economy}, 88(3), 495--503.
\item Limão, N. (2006). Preferential Trade Agreements as Stumbling Blocks for Multilateral Trade Liberalization: Evidence for the United States. \emph{American Economic Review}, 96(3), 896--914.
\item McCallum, J. (1995). National Borders Matter: Canada-U.S. Regional Trade Patterns. \emph{American Economic Review}, 85(3), 615--623.
\item Meese, R. y K. Rogoff (1983). Empirical Exchange Rate Models of the Seventies: Do They Fit Out of Sample? \emph{Journal of International Economics}, 14, 3--24.
\item Melitz, M. (2003). The Impact of Trade on Intra-Industry Reallocations and Aggregate Industry Productivity. \emph{Econometrica}, 71(6), 1695--1725.
\item Mendoza, E. (1991). Real Business Cycles in a Small Open Economy. \emph{American Economic Review}, 81(4), 797--818.
\item Mendoza, E. (2010). Sudden Stops, Financial Crises, and Leverage. \emph{American Economic Review}, 100(5), 1941--1966.
\item Mongelli, F. (2002). «New» Views on the Optimum Currency Area Theory: What Is EMU Telling Us? ECB Working Paper 138.
\item Neary, J. P. (1978). Short-Run Capital Specificity and the Pure Theory of International Trade. \emph{Economic Journal}, 88(351), 488--510.
\item Obstfeld, M. \emph{The Logic of Currency Crises}. Universidad de California en Berkeley.
\item Obstfeld, M. y K. Rogoff (1995). The Mirage of Fixed Exchange Rates. \emph{Journal of Economic Perspectives}, 9(4), 73--96.
\item Obstfeld, M. y K. Rogoff (1996). \emph{Foundations of International Macroeconomics}. The MIT Press.
\item Obstfeld, M. y K. Rogoff (2000). The Six Major Puzzles in International Macroeconomics: Is There a Common Cause? \emph{NBER Macroeconomics Annual}, 15, 339--390.
\item Ostry, J., A. Ghosh, M. Chamon y M. Qureshi (2012). Tools for Managing Financial-Stability Risks from Capital Inflows. \emph{Journal of International Economics}, 88(2), 407--421.
\item Panagariya, A. (2000). Preferential Trade Liberalization: The Traditional Theory and New Developments. \emph{Journal of Economic Literature}, 38(2), 287--331.
\item Pavcnik, N. (2017). The Impact of Trade on Inequality in Developing Countries. NBER Working Paper 23878.
\item Rodríguez, G., P. Villanueva y P. Castillo (2018). Driving Economic Fluctuations in Peru: The Role of the Terms of Trade. \emph{Empirical Economics}, 55(3), 1089--1119.
\item Rose, A. (2000). One Money, One Market: Estimating the Effect of Common Currencies on Trade. \emph{Economic Policy}, 15(30), 7--45.
\item Rybczynski, T. M. (1955). Factor Endowment and Relative Commodity Prices. \emph{Economica}, 22(88), 336--341.
\item Schmitt-Grohé, S. y M. Uribe (2003). Closing Small Open Economy Models. \emph{Journal of International Economics}, 61, 163--185.
\item Schmitt-Grohé, S. y M. Uribe (2018). How Important Are Terms-of-Trade Shocks? \emph{International Economic Review}, 59(1), 85--111.
\item Stolper, W. y P. Samuelson (1941). Protection and Real Wages. \emph{Review of Economic Studies}, 9(1), 58--73.
\item Taylor, A. y M. Taylor (2004). The Purchasing Power Parity Debate. \emph{Journal of Economic Perspectives}, 18(4), 135--158.
\item Tovar, P. (2012). Preferential Trade Agreements and Unilateral Liberalization: Evidence from CAFTA. \emph{World Trade Review}, 11(4), 591--619.
\item Tovar, P. (2019). Preferential and Multilateral Liberalization: Evidence from Latin America's Use of Tariffs, Antidumping and Safeguards. \emph{Journal of Development Economics}, 141.
\item Trefler, D. (1993). International Factor Price Differences: Leontief Was Right! \emph{Journal of Political Economy}, 101(6), 961--987.
\item Trefler, D. (1995). The Case of the Missing Trade and Other Mysteries. \emph{American Economic Review}, 85(5), 1029--1046.
\end{itemize}
}

\subsection*{E.4 Otras referencias}
{\small
\begin{itemize}[leftmargin=1.4em,itemsep=1pt,label={}]
\item Castillo, P., C. Montoro y V. Tuesta (2009). Un modelo de equilibrio general con dolarización para la economía peruana. \emph{Revista Estudios Económicos}, BCRP, 17, 9--50.
\item Gopinath, G. (2019). A Case for an Integrated Policy Framework. Jackson Hole Economic Policy Symposium, Federal Reserve Bank of Kansas City.
\item Montoro, C. y M. Ortiz (2016). Foreign Exchange Intervention and Monetary Policy Design: A Market Microstructure Analysis. BCRP, Documento de Trabajo.
\end{itemize}
}

\vspace{1.2cm}
\franjas[2pt]
\begin{center}\small\color{tonoPizarra}
Plan de lectura de elaboración propia. Los títulos de los capítulos
corresponden a las obras citadas y se reproducen únicamente con fines de
referencia bibliográfica. La asignación de capítulos y secciones se basa en
los índices de los libros de la carpeta; los temas de los cursos proceden de
los programas citados en el Anexo C.
\end{center}

\end{document}
